\documentclass[11pt]{article}

% ----- Codificacion y fuentes (xelatex) -----
\usepackage{fontspec}
\usepackage{polyglossia}
\setmainlanguage{spanish}
\setmainfont{Calibri}[
  BoldFont      = Calibri Bold,
  ItalicFont    = Calibri Italic,
  BoldItalicFont= Calibri Bold Italic ]
\newfontfamily\headingfont{Calibri}

% ----- Geometria y espaciado -----
\usepackage[a4paper,margin=2.3cm]{geometry}
\usepackage{setspace}
\setstretch{1.12}
\usepackage{parskip}

% ----- Color y cajas -----
\usepackage[table]{xcolor}
\definecolor{agua}{HTML}{0F8B8D}
\definecolor{aguaOsc}{HTML}{04302F}
\definecolor{verde}{HTML}{2FB37F}
\definecolor{ambar}{HTML}{E0A800}
\definecolor{rojo}{HTML}{D64545}
\definecolor{grisc}{HTML}{F0F5F5}
\definecolor{lineac}{HTML}{CBD8D8}

\usepackage[most]{tcolorbox}
\newtcolorbox{aviso}{colback=rojo!8,colframe=rojo!70,boxrule=0.8pt,
  arc=2pt,left=8pt,right=8pt,top=6pt,bottom=6pt}
\newtcolorbox{nota}{colback=agua!8,colframe=agua!70,boxrule=0.8pt,
  arc=2pt,left=8pt,right=8pt,top=6pt,bottom=6pt}

% ----- Tablas -----
\usepackage{booktabs}
\usepackage{tabularx}
\usepackage{array}
\renewcommand{\arraystretch}{1.35}
\newcolumntype{L}{>{\raggedright\arraybackslash}X}

% ----- Titulos de seccion -----
\usepackage{titlesec}
\titleformat{\section}{\headingfont\Large\bfseries\color{aguaOsc}}{\thesection.}{0.5em}{}
\titleformat{\subsection}{\headingfont\large\bfseries\color{agua}}{\thesubsection}{0.5em}{}
\titlespacing*{\section}{0pt}{16pt}{6pt}
\titlespacing*{\subsection}{0pt}{10pt}{4pt}

% ----- Listas -----
\usepackage{enumitem}
\setlist[itemize]{leftmargin=1.4em,itemsep=2pt,topsep=3pt}

% ----- Encabezado/pie -----
\usepackage{fancyhdr}
\pagestyle{fancy}
\fancyhf{}
\renewcommand{\headrulewidth}{0.4pt}
\renewcommand{\footrulewidth}{0.4pt}
\fancyhead[L]{\small\color{agua}Alerta temprana de biomasa algal (HABs)}
\fancyhead[R]{\small\color{agua}Manual de usuario}
\fancyfoot[C]{\small\thepage}

\usepackage{hyperref}
\hypersetup{colorlinks=true,linkcolor=agua,urlcolor=agua,
  pdftitle={Manual de usuario - Alerta temprana de biomasa algal (HABs)}}

% Etiquetas de color para reemplazar emojis
\newcommand{\lbl}[2]{\textcolor{#1}{\textbf{[#2]}}}

\begin{document}

% ============ PORTADA ============
\begin{titlepage}
\centering
\vspace*{2.2cm}
{\color{lineac}\rule{\textwidth}{1.2pt}}\\[1.2cm]
{\headingfont\Huge\bfseries\color{aguaOsc} Manual de Usuario}\\[0.9cm]
{\headingfont\LARGE\color{agua} Sistema de Alerta Temprana\\[0.3cm] de Biomasa Algal (HABs)}\\[0.8cm]
{\Large Pron\'ostico de clorofila-a a 0--7 d\'ias}\\[1.2cm]
{\color{lineac}\rule{\textwidth}{1.2pt}}\\[2.5cm]

\begin{minipage}{0.85\textwidth}\centering\itshape\color{aguaOsc}
Herramienta de apoyo a la decisi\'on basada en im\'agenes satelitales Sentinel-2
y modelos de aprendizaje autom\'atico. No sustituye la verificaci\'on de campo.
\end{minipage}

\vfill
{\large Interfaz web: \texttt{streamlit run app.py}}\\[0.4cm]
\end{titlepage}

\tableofcontents
\newpage

% ============ 1 ============
\section{\'Qu\'e es este programa?}

Es una herramienta de \textbf{alerta temprana} que \textbf{pronostica} la biomasa
algal ---usando la \textbf{clorofila-a como proxy}--- en cuerpos de agua, con
\textbf{0 a 7 d\'ias de anticipaci\'on}.

Trabaja a partir de \textbf{im\'agenes satelitales Sentinel-2} y modelos ya
entrenados. No reentrena nada: envuelve la l\'ogica de pron\'ostico existente en una
interfaz web sencilla (Streamlit).

\begin{aviso}
\textbf{Muy importante:} es un \textbf{PRON\'OSTICO a futuro}, no una detecci\'on
sobre la imagen que usted sube. Estima c\'omo estar\'a la biomasa dentro de 1, 3, 5
o 7 d\'ias.
\end{aviso}

\subsection*{Lo que S\'I hace}
\begin{itemize}
  \item Estima la \textbf{clorofila-a media prevista} (\textmu g/L) a +1, +3, +5 y +7 d\'ias.
  \item Genera un \textbf{mapa de biomasa algal prevista} por p\'ixel.
  \item Da una \textbf{alerta por nivel} (Normal / Elevada / Floraci\'on).
  \item Calcula una \textbf{banda de incertidumbre} (P10--P90, \textasciitilde80\% de confianza).
  \item Estima la \textbf{probabilidad de una anomal\'ia} (salto at\'ipico para ese cuerpo).
  \item Explica \textbf{por qu\'e} (variables m\'as influyentes, SHAP).
  \item Permite \textbf{descargar} mapa, animaci\'on y resumen CSV.
\end{itemize}

\subsection*{Lo que NO hace}
\begin{itemize}
  \item \textbf{No confirma toxicidad} ni identifica la especie (cianobacterias, marea
  roja). La clorofila-a mide el \textbf{pigmento}, no la toxina ni el organismo.
  \item \textbf{No acepta fotos comunes} (RGB de celular ni capturas de Google Maps).
  \item \textbf{No funciona fuera de los 5 cuerpos validados.}
\end{itemize}

% ============ 2 ============
\section{Requisitos}
\begin{itemize}
  \item \textbf{Python} con las dependencias del proyecto instaladas.
  \item \textbf{Modelos de producci\'on} presentes en \texttt{artifacts/models/}
  (si faltan, se generan con \texttt{python train\_final.py}).
  \item Un navegador web (la app abre en el navegador autom\'aticamente).
\end{itemize}

% ============ 3 ============
\section{C\'omo abrir el programa}
Desde una terminal, ubicado en la carpeta \texttt{habs\_forecast}:
\begin{quote}\ttfamily streamlit run app.py\end{quote}
Se abrir\'a una p\'agina en el navegador con el t\'itulo
\textbf{``Alerta temprana de biomasa algal (HABs)''}.

% ============ 4 ============
\section{Cuerpos de agua disponibles (validados)}
\begin{tabularx}{\textwidth}{L L L}
\toprule
\textbf{Cuerpo} & \textbf{Tipo} & \textbf{Pa\'is} \\
\midrule
Lago Okeechobee   & Lago / agua dulce        & Estados Unidos \\
Bah\'ia de Tampa   & Costa / marino-estuarino & Estados Unidos \\
Embalse El Caj\'on & Lago / agua dulce        & Honduras \\
Golfo de Fonseca  & Costa / marino-estuarino & Honduras \\
Lago de Yojoa     & Lago / agua dulce        & Honduras \\
\bottomrule
\end{tabularx}

\begin{nota}
Algunos cuerpos pueden marcarse como \textbf{EXPLORATORIOS}: no tienen verdad de campo
in-situ en la ventana 2023--2026 y tienen menos datos. Sus resultados son de
\textbf{menor confianza} y la app lo advierte en pantalla.
\end{nota}

% ============ 5 ============
\section{Recorrido por la interfaz}

\subsection{Barra lateral (izquierda) --- ``C\'omo leer esta herramienta''}
Explica las reglas clave del sistema (es pron\'ostico, requiere Sentinel-2 de 5 bandas,
solo 5 cuerpos, clorofila = proxy). L\'eala una vez antes de usar la herramienta.

\subsection{Selectores (parte superior)}
\begin{enumerate}[leftmargin=1.6em]
  \item \textbf{Cuerpo de agua} --- elija uno de los 5 cuerpos validados.
  \item \textbf{Horizonte de pron\'ostico} --- +1, +3, +5 o +7 d\'ias (por defecto \textbf{+3}).
  \begin{itemize}
    \item \textbf{+1 y +7 d\'ias} son horizontes \textit{body-level}: el modelo predice el
    \textbf{nivel del cuerpo completo} y el mapa reparte ese nivel seg\'un el patr\'on
    espacial actual (es una estimaci\'on, no un pron\'ostico p\'ixel a p\'ixel).
    \item \textbf{+3 y +5 d\'ias} usan se\~nal espectral por p\'ixel: el mapa muestra un
    \textbf{gradiente real} con detalle espacial.
  \end{itemize}
  \item \textbf{Tipo / Pa\'is} --- informativo (se completa solo seg\'un el cuerpo elegido).
\end{enumerate}

\subsection{Entrada de escena Sentinel-2}
\textbf{A) Usar escena de ejemplo} (recomendado para la demostraci\'on)
\begin{itemize}
  \item La app \textbf{ordena las escenas por calidad} (agua m\'as limpia primero).
  \item Con la casilla \textit{``Usar autom\'aticamente la mejor escena''} marcada, elige
  sola la mejor fecha (evita escenas nubladas donde el agua ni se ve).
  \item Si la desmarca, puede elegir la fecha manualmente en el desplegable (la mejor
  aparece marcada como mejor).
  \item Si la escena elegida tiene nubosidad/neblina, aparece una advertencia.
\end{itemize}

\textbf{B) Subir GeoTIFF}
\begin{itemize}
  \item Debe ser un \textbf{raster georreferenciado Sentinel-2 de 5 bandas} en el orden
  \textbf{B2 (azul), B3 (verde), B4 (rojo), B5 (red-edge), B8 (NIR)}.
  \item Las bandas red-edge (B5) e infrarrojo (B8) son las que estiman la clorofila; por
  eso \textbf{una foto RGB com\'un es rechazada}.
  \item El contexto no espectral (clorofila reciente, ERA5, in-situ) se toma de la \'ultima
  fecha disponible del cuerpo.
\end{itemize}

\subsection{Bot\'on ``Analizar''}
\begin{itemize}
  \item (Opcional) Marque \textbf{``Mostrar animaci\'on tipo pron\'ostico del clima''} para
  generar un video corto que recorre +1 $\rightarrow$ +7 d\'ias. Debe marcarse
  \textbf{antes} de Analizar.
  \item Pulse \textbf{Analizar}. Aparece un indicador de progreso mientras procesa la
  escena y genera el pron\'ostico.
\end{itemize}

% ============ 6 ============
\section{C\'omo leer los resultados}
Tras ``Analizar'' se muestra, de arriba hacia abajo:

\subsection{Encabezado}
Cuerpo, horizonte, fecha de la escena (\textbf{t0}) y \textbf{etiqueta de confianza}
(seg\'un frescura de datos, cobertura de p\'ixeles de agua y estado).

\subsection{Mapa (2 paneles)}
\begin{itemize}
  \item \textbf{Izquierda:} imagen satelital real de la escena.
  \item \textbf{Derecha:} \textbf{biomasa algal prevista} por p\'ixel para el horizonte elegido.
\end{itemize}

\subsection{Insignia de credibilidad}
Compara el \textit{skill} validado del modelo frente a la persistencia (l\'inea base) en
ese horizonte.

\subsection{Animaci\'on (si la activ\'o)}
Video en bucle que arranca en el estado observado de hoy y recorre la biomasa prevista a
1, 3, 5 y 7 d\'ias, como un pron\'ostico del clima en la tele.

\subsection{Nivel de biomasa algal (banner de color)}
\begin{tabularx}{\textwidth}{c >{\bfseries}l L}
\toprule
\textbf{Color} & \textbf{Nivel} & \textbf{Significado} \\
\midrule
\lbl{verde}{VERDE}  & NORMAL          & Clorofila prevista por debajo del umbral elevado \\
\lbl{ambar}{\'AMBAR} & BIOMASA ELEVADA & Clorofila prevista por encima del umbral elevado \\
\lbl{rojo}{ROJO}    & FLORACI\'ON      & Clorofila prevista igual o mayor al umbral de floraci\'on \\
\bottomrule
\end{tabularx}

Debajo se indica el \textbf{\% de \'area} en floraci\'on y en biomasa elevada.

\subsection{Dos se\~nales que pueden diferir (\textit{no confundir})}
\begin{itemize}
  \item \textbf{NIVEL (magnitud):} cu\'anta clorofila se prev\'e, comparada con umbrales
  biol\'ogicos (es el banner de color).
  \item \textbf{Probabilidad de anomal\'ia (P85):} la probabilidad de un \textbf{salto
  at\'ipico para ESE cuerpo} (clasificador calibrado).
\end{itemize}
Un embalse cr\'onicamente alto puede dar \textbf{nivel alto} y \textbf{anomal\'ia baja}:
ese nivel es \textit{su} normal, no un evento inusual. Y al rev\'es. La app muestra una
nota aclaratoria cuando ambas se\~nales divergen.

\subsection{Clorofila-a prevista (intensidad)}
Valor medio previsto en \textmu g/L con su \textbf{banda de incertidumbre P10--P90}
(calibrada, CQR \textasciitilde80\%).

\subsection{Pesta\~nas din\'amicas}
\begin{itemize}
  \item \textbf{Trayectoria 0--7 d\'ias:} clorofila prevista en cada horizonte con su banda;
  el punto grande es el horizonte seleccionado, el color indica el nivel.
  \item \textbf{Medidor de riesgo:} probabilidad calibrada de salto an\'omalo a +N d\'ias.
  La l\'inea roja es el \textbf{umbral operativo real}: por encima, el sistema dispara
  alerta. El umbral es bajo a prop\'osito (prioriza no perder eventos $\rightarrow$ alta
  sensibilidad).
  \item \textbf{\'Por qu\'e? (SHAP):} variables que m\'as pesan en el pron\'ostico. En corto
  plazo domina la clorofila reciente; a mayor horizonte entran meteorolog\'ia, nutrientes
  e \'indices espectrales.
\end{itemize}

% ============ 7 ============
\section{Descargas}
En la secci\'on \textbf{Descargas} hay tres botones:

\begin{tabularx}{\textwidth}{>{\bfseries}l l L}
\toprule
\textbf{Bot\'on} & \textbf{Archivo} & \textbf{Contenido} \\
\midrule
Mapa (PNG)       & \texttt{mapa\_<cuerpo>\_h<N>\_<fecha>.png} & Los dos paneles del mapa \\
Animaci\'on (GIF) & \texttt{animacion\_<cuerpo>\_<fecha>.gif}   & Solo si activ\'o la animaci\'on \\
Pron\'ostico (CSV)& \texttt{pronostico\_<cuerpo>\_<fecha>.csv}  & Todos los horizontes (0--7 d) \\
\bottomrule
\end{tabularx}

El \textbf{CSV} incluye, por horizonte: clorofila prevista, P10, P90, probabilidad de
alerta y nivel; m\'as una cabecera con metadatos y el recordatorio de que \textbf{no
confirma toxicidad}.

% ============ 8 ============
\section{Interpretaci\'on responsable}
\begin{itemize}
  \item La clorofila-a es un \textbf{proxy de biomasa}, \textbf{no} una medida de toxicidad
  ni de la especie. Una alerta indica \textbf{riesgo que amerita verificaci\'on de campo}.
  \item Confirmar cianobacterias, dinoflagelados o toxinas requiere \textbf{muestreo de
  campo} (microscop\'ia / ensayos de toxinas).
  \item Fuera de los 5 cuerpos validados \textbf{no hay modelo ni calibraci\'on}: no use la
  herramienta ah\'i.
  \item Trate los resultados de cuerpos \textbf{exploratorios} con cautela.
\end{itemize}

\begin{aviso}
\textbf{Proxy de biomasa algal (clorofila-a).} NO confirma toxicidad ni floraci\'on
nociva. Herramienta de \textbf{alerta temprana}; requiere \textbf{verificaci\'on de campo}.
\end{aviso}

% ============ 9 ============
\section{Problemas frecuentes}
\begin{tabularx}{\textwidth}{L L L}
\toprule
\textbf{Situaci\'on} & \textbf{Causa probable} & \textbf{Qu\'e hacer} \\
\midrule
``Faltan los modelos de producci\'on'' & No est\'an en \texttt{artifacts/models/} & Ejecutar \texttt{python train\_final.py} \\
``No hay modelo entrenado para\ldots a +N d\'ias'' & Falta ese modelo/horizonte & Entrenar o elegir otro horizonte \\
``El archivo NO tiene 5 bandas v\'alidas'' & Subi\'o una foto RGB o un TIFF incompleto & Usar un GeoTIFF Sentinel-2 de 5 bandas \\
``Muy pocos p\'ixeles de agua v\'alidos'' & Escena nublada o sin agua & Probar otra fecha o la mejor escena autom\'atica \\
Advertencia de nubosidad/neblina & La escena tiene poco agua clara & Usar la mejor escena autom\'atica u otra fecha \\
``Explicabilidad SHAP no disponible'' & No se gener\'o SHAP & Ejecutar \texttt{python explain\_model.py} \\
El bot\'on ``Analizar'' est\'a gris & No hay escena cargada & Seleccionar escena de ejemplo o subir GeoTIFF v\'alido \\
\bottomrule
\end{tabularx}

% ============ 10 ============
\section{Ficha t\'ecnica (resumen)}
\begin{itemize}
  \item \textbf{Modelo de intensidad + intervalos:} XGBoost con CQR (regresi\'on cuant\'ilica
  conformalizada).
  \item \textbf{Alerta:} red neuronal (HABNet), por grupo ecol\'ogico y horizonte.
  \item \textbf{Validaci\'on:} pron\'ostico causal sin fuga de informaci\'on (validaci\'on
  anidada).
  \item \textbf{Entrada:} escena Sentinel-2 de 5 bandas (B2, B3, B4, B5, B8) + contexto no
  espectral (clorofila reciente, ERA5, in-situ).
  \item \textbf{Salida:} clorofila-a prevista, banda P10--P90, nivel, probabilidad de
  anomal\'ia, mapa por p\'ixel y explicaci\'on SHAP.
\end{itemize}

\vspace{1cm}
{\color{lineac}\rule{\textwidth}{0.6pt}}\\[0.3cm]
{\small\itshape\color{aguaOsc} Documento generado como manual de usuario del sistema de
alerta temprana de biomasa algal (HABs). Herramienta de apoyo a la decisi\'on; no
sustituye la verificaci\'on de campo.}

\end{document}
